\subsection{স্থানমান, বৈধ অঙ্ক ও সংখ্যার পরিসর}
\textbf{স্থানমান} নির্ধারিত হয় অঙ্কের অবস্থান ও ভিত্তি দিয়ে। ভিত্তি $b$ হলে বিন্দুর বামে স্থানগুলোর ওজন $b^0,b^1,b^2,\ldots$ এবং ডানে $b^{-1},b^{-2},\ldots$। পূর্ণাংশ ও ভগ্নাংশের বিভাজককে \textbf{র‍্যাডিক্স পয়েন্ট} বলে; দশমিক পদ্ধতিতে এটি দশমিক বিন্দু। অঙ্কের নিজস্ব মান এবং অঙ্কটির অবদান আলাদা: $(343)_5$-এ ডানদিকের ৩-এর অবদান ৩, বামদিকের ৩-এর অবদান $3\times25=75$।

ভিত্তি $b$-এর প্রতিটি অঙ্কের মান অবশ্যই $0\leq d<b$ হবে। তাই $(128)_8$ অবৈধ; অক্টালে ৮ নেই। $(1G)_{16}$-ও অবৈধ; হেক্সাডেসিমালে শেষ প্রতীক F। কোনো লেখায় ব্যবহৃত বৃহত্তম অঙ্ক ৭ হলে তার ভিত্তি \emph{অন্তত} ৮; শুধু অঙ্ক দেখে ভিত্তি ঠিক ৮ বলা যায় না।

\chThreeExample{১}{অজানা ভিত্তি নির্ণয়}
$(132)_b=(56)_{10}$ হলে $b$ কত?
\chThreeAnswer
\[b^2+3b+2=56\Rightarrow b^2+3b-54=0\Rightarrow(b+9)(b-6)=0.\]
গ্রহণযোগ্য ভিত্তি $b=6$; $-9$ গ্রহণযোগ্য নয়। যাচাই: $36+18+2=56$।

\chThreeExample{২}{যোগের সমতা থেকে ভিত্তি}
$(12)_b+(15)_b=(30)_b$।
\chThreeAnswer
\[(b+2)+(b+5)=3b\Rightarrow2b+7=3b\Rightarrow b=7.\]
বৃহত্তম অঙ্ক ৫ ভিত্তি ৭-এর চেয়ে ছোট। দশমিকে $9+12=21$; সাত ভিত্তিতে $21=(30)_7$।

\chThreeExample{৩}{একই ১০-এর মান কেন বদলায়?}
বাইনারিতে $1+1=(10)_2$, অক্টালে $7+1=(10)_8$, হেক্সাডেসিমালে $F+1=(10)_{16}$। প্রতিটি ক্ষেত্রে ডানদিকের ঘরের অঙ্ক শেষ হলে বাম ঘরে ১ ক্যারি যায়। তাই $(10)_b$-এর দশমিক মান $b$; ১০ লেখা দেখলেই দশ বলা যাবে না।

\smallhead{বিট, বাইট ও ধারণক্ষমতা}
একটি বাইনারি অঙ্ক হলো \textbf{বিট}; ৮ বিটে এক বাইট এবং ৪ বিটে এক নিবল। প্রসেসর একবারে যে দৈর্ঘ্যের তথ্য নিয়ে কাজ করার জন্য তৈরি, তাকে তার \textbf{ওয়ার্ড দৈর্ঘ্য} বলা হয়; ওয়ার্ড সব প্রসেসরে এক বাইট নয়। $n$ বিটে $2^n$টি pattern হয়। চিহ্নবিহীন পূর্ণসংখ্যা হিসেবে সর্বনিম্ন মান ০ এবং সর্বোচ্চ মান $2^n-1$।

\chThreeExample{৪}{অবস্থার সংখ্যা ও সর্বোচ্চ মান}
৫ বিটে কতটি পৃথক অবস্থা ও সর্বোচ্চ কত মান পাওয়া যায়?
\chThreeAnswer $2^5=32$টি অবস্থা; চিহ্নবিহীন সর্বোচ্চ মান $31=(11111)_2$। ৩২টি অবস্থা মানে সর্বোচ্চ মান ৩২ নয়, কারণ ০-ও একটি অবস্থা।

\chThreeExample{৫}{ধারণাভিত্তিক MCQ}
অক্টাল সংখ্যার একটি অঙ্ক প্রকাশে কতটি বাইনারি বিট দরকার? (ক) ২, (খ) ৩, (গ) ৪, (ঘ) ৮।
\chThreeAnswer (খ) ৩। কারণ $2^2<8=2^3$। কিন্তু একটি অক্টাল অঙ্ক এবং পুরো অক্টাল সংখ্যা এক বিষয় নয়।

% Mahbub teaching-type follow-up, 2026-10-14
\chThreeExample{৬}{আগের ও পরের সংখ্যা}
$(77)_8$, $(2F)_{16}$ ও $(1111)_2$-এর ঠিক আগের ও পরের পূর্ণসংখ্যা লেখো।
\chThreeAnswer
\begin{center}\begin{tabular}{|c|c|c|}\hline
আগের সংখ্যা & প্রদত্ত সংখ্যা & পরের সংখ্যা\\\hline
$(76)_8$ & $(77)_8$ & $(100)_8$\\\hline
$(2E)_{16}$ & $(2F)_{16}$ & $(30)_{16}$\\\hline
$(1110)_2$ & $(1111)_2$ & $(10000)_2$\\\hline
\end{tabular}\end{center}
অক্টালে $7+1=10_8$; ফলে 77-এর ডানের ঘরের ক্যারি বামের ঘরেও যায়। হেক্সাডেসিমালে F-এর পরে ০ রেখে ক্যারি ১ দিতে হয়। নির্দিষ্ট চার বিটের register-এ 1111-এর পরে carry বাদ দিলে 0000; কিন্তু বিটসংখ্যা সীমাবদ্ধ না করলে পূর্ণসংখ্যার পরের মান 10000।

\chThreeExample{৭}{অঙ্কের স্থান বদলালে মানের পরিবর্তন}
$(35)_8$ ও $(53)_8$-এর পার্থক্য কত? ৫ কোনটিতে বেশি অবদান রাখে?
\chThreeAnswer $(35)_8=3\times8+5=29$; $(53)_8=5\times8+3=43$। পার্থক্য $14=(16)_8$। প্রথমটিতে ৫-এর অবদান ৫, দ্বিতীয়টিতে ৪০। অঙ্ক একই থাকলেও position বদলালে সংখ্যার মান বদলায়।
